\begin{claim}{B.16u}{Proved}
Let $(\mathcal C,\mathcal D)/R$ be a stable marked model in the
residue-characteristic-zero setting of Part B. Put
$J=\rho(I_K)\subseteq\Out(\Pi)$, with the full geometric profinite
group $\Pi$. The pair $(\Pi,J)$ determines the set
\[
 \{([\Pi_e]_{\Pi},m_e):e\text{ a geometric node}\}.
\]
This holds for affine curves without a supplied set of cusp classes,
and is functorial under every continuous isomorphism of the geometric
groups carrying one closed inertia image onto the other. No generator
of $J$ is part of the data. The thicknesses use the normalized valuation
of the trait over which the stable model is given.
\end{claim}
\begin{proof}
An unramified finite extension splitting the special fibre changes neither
geometric groups, geometric node classes, thicknesses nor inertia. We can
therefore split the fibre for the proof without adding data to the
reconstruction. If $J=1$, B.5 says there are no nodes. Suppose there is a
node. Then B.5 identifies $J$ with $\widehat\Z$ and, after a complex
topological comparison and a choice of orientation, identifies a
topological generator with the completion of
\[
 \varphi=\prod_e\tau_e^{m_e}\in\operatorname{Mod}(S),
\]
where $S$ is a compact oriented surface whose boundary circles represent
the cusps. The node circles form a disjoint essential multicurve, and
every component $S_v$ cut out by that multicurve has $\chi(S_v)<0$.
These facts are exactly stability: a disk or an annulus component would
be an unstable rational component, and a peripheral or parallel pair
of node circles would cut off such a component. A base point outside
the twist annuli gives an actual automorphism of
$\Gamma=\pi_1(S)$. We use this choice only to prove the reconstruction
formula; it is not retained as input.

\paragraph{The group constructed from the inertia image.}
Because $\Pi$ is finitely generated and centre-free, B.2 constructs,
from the given pair alone, the profinite extension
\[
 1\longrightarrow\Pi\longrightarrow
 E_J:=\Aut(\Pi)\times_{\Out(\Pi)}J
 \stackrel{q}{\longrightarrow}J\longrightarrow1.
\]
The chosen automorphism $\widehat\varphi$ generates a procyclic
subgroup of $\Aut(\Pi)$ which maps isomorphically onto $J$:
the composite $\widehat\Z\to\Aut(\Pi)\to J$ is an isomorphism
by B.5. Consequently
\[
 E_J\simeq\Pi\rtimes_{\widehat\varphi}\widehat\Z
       \simeq\widehat{\Gamma\rtimes_{\varphi}\Z}.
\]
For completeness, the second isomorphism is full profinite completion,
not a procyclic approximation. Characteristic finite-index subgroups
of the finitely generated group $\Gamma$ are cofinal among its
finite-index subgroups. Thus $\varphi$ extends to the indicated
continuous action on $\widehat\Gamma=\Pi$. Every homomorphism
$\Gamma\rtimes\Z\to Q$ to a finite group extends on $\Gamma$ to
$\Pi$ and on $\Z$ to $\widehat\Z$; the semidirect relation extends
by continuity. Conversely their product has dense image. These facts
give both universal maps and show they are inverse. The marked normal
subgroup $\Pi\subset E_J$ is preserved throughout.

\paragraph{The mapping torus and its geometric decomposition.}
Let $M_\varphi$ be the mapping torus of the multitwist, represented by a
homeomorphism equal to the identity off the disjoint twist annuli and on
$\partial S$. It is compact and orientable, with one boundary torus
for each cusp, and its fundamental group is
$\Gamma\rtimes_\varphi\Z$. The interior has universal cover
$\widetilde{\operatorname{int}S}\times\R\simeq\R^3$.
Irreducibility descends from this cover by
\citet[Proposition 1.6]{BC_Hatcher3M}. Each boundary torus is
incompressible: in the surface-bundle extension its group is generated
by a nontrivial boundary loop and a lift of the generator of $\Z$,
and projection to $\Z$ first kills the second exponent in any proposed
relation, after which surface-group injectivity kills the first.
The same argument applies to each node torus.

Cutting at the node tori gives the product Seifert pieces
$M_v=S_v\times S^1$. There are no solid-torus, annular-product or
twisted-Klein-bottle pieces since $\chi(S_v)<0$. These tori are exactly
the JSJ tori. Here are the relevant minimality checks. None is
boundary-parallel, and no two are parallel: either possibility would
make the corresponding essential surface circles peripheral or
freely homotopic, respectively. Such circles are excluded by the stable
multicurve. On the torus associated to $e$, choose the surface circle
$c_e$ and the circle fibre $t_v$ on one side. The circle fibre on the
other side is
\[
             t_w=c_e^{\epsilon_e m_e}t_v,
             \qquad \epsilon_e\in\{1,-1\}.
\tag{B.JSJ.1}
\]
Indeed the two product trivialisations of the annulus mapping torus
differ by exactly $m_e$ full twists. This also follows directly from
the separating-edge or HNN formula in B.5. Thus the two primitive
fibre slopes have intersection number $m_e>0$.

Deleting any of these cuts leaves an essential torus, so the joined
piece is not atoroidal. It cannot be Seifert either. In an irreducible
Seifert piece an essential torus is vertical or horizontal
\citep[Proposition 1.12]{BC_Hatcher3M}. A horizontal torus gives a
finite torus cover of the base orbifold and hence a virtually solvable
fundamental group, whereas the joined piece contains the nonabelian
free group $\pi_1(S_v)$. It must therefore be vertical. The induced
Seifert fibrations on each $S_v\times S^1$ have the product fibre
slope on the boundary: the only exceptions to boundary-slope uniqueness
are a solid torus or an $I$-bundle over a torus or Klein bottle
\citep[Lemma 1.17]{BC_Hatcher3M}, already excluded here.
Their slopes would have to coincide across the omitted cut, contrary
to (B.JSJ.1). This argument also applies to several omitted cuts and
to a loop edge. Minimality and uniqueness in the torus-decomposition
theorem identify the displayed decomposition with the JSJ decomposition
\citep[Theorem 1.9]{BC_Hatcher3M}.

\paragraph{The exact profinite JSJ input.}
We use \citet[Theorem 5.5]{BC_Xu2025}, in the opened arXiv version~4,
whose hypotheses are compact orientable irreducible $3$-manifolds with
empty or incompressible toral boundary, excluding closed Sol manifolds.
It states that every isomorphism of their profinite fundamental groups
induces a congruent isomorphism of their JSJ graphs of profinite groups.
In particular it preserves vertex and edge groups up to ambient
conjugacy, with the two incidence homomorphisms. No peripheral marking
is a hypothesis. The same source, Propositions 5.2--5.4, supplies
efficiency, the standard profinite tree and the boundary version used
in that theorem. Our manifold satisfies these hypotheses; when there
are nodes its nonabelian product pieces exclude the Sol exception.
This is a different input from the cuspidal or noncuspidal hypotheses
of the combinatorial anabelian results used in B.16.

Accordingly $E_J$ recovers the profinite groups
\[
 G_v=\widehat{\pi_1(S_v)}\times\widehat\Z,
 \qquad A_e=\widehat{\pi_1(T_e)}\simeq\widehat\Z^2,
\]
and their incidences. The geometric node group is recovered by
\[
                         \Pi_e=A_e\cap\Pi.
\tag{B.JSJ.2}
\]
This is an equality with the full closed cyclic subgroup: the surface
bundle map on $A_e=\overline{\langle c_e,t_v\rangle}$ sends $c_e$
to zero and $t_v$ to a generator of $J$, so its kernel is precisely
$\overline{\langle c_e\rangle}$.

The JSJ theorem initially gives ambient $E_J$ conjugacy. This yields
exactly the required $\Pi$ conjugacy. For any $g\in E_J$, choose
$h\in A_e$ with $q(h)=q(g)$; such an $h$ exists because
$q(A_e)=J$. Then $gh^{-1}\in\Pi$, and $h$ normalizes
$A_e\cap\Pi$. Hence the $E_J$ conjugates of this node group are
already its $\Pi$ conjugates. The same adjustment applies after an
isomorphism between two reconstructed extensions.

\paragraph{An intrinsic integral thickness formula.}
The surface factor of $G_v$ is a free profinite group of rank at least
two, hence centre-free; thus $Z(G_v)$ is exactly its product
$\widehat\Z$ fibre. For the two incidence embeddings
$\iota_{e,i}:A_e\hookrightarrow G_{v_i}$, put
\[
 F_{e,i}=\iota_{e,i}^{-1}(Z(G_{v_i})),\qquad i=0,1.
\]
Both are closed primitive cyclic subgroups of $A_e$. They remain two
different subgroups when $v_0=v_1$, because the two incidence embeddings
are retained. Formula (B.JSJ.1), completed profinitely, gives
\[
 \boxed{\quad m_e=[A_e:F_{e,0}F_{e,1}].\quad}
\tag{B.JSJ.3}
\]
In the basis $(c_e,t_v)$ the two columns are $(0,1)$ and
$(\epsilon_em_e,1)$, so their cokernel is
$\widehat\Z/m_e\widehat\Z$, of order $m_e$.
This finite index is preserved by every congruent isomorphism of the
JSJ groups. It uses no orientation or chosen generator of $J$.
It also explains directly why retaining all primes recovers the full
positive integer rather than only selected valuations.

Finally an isomorphism $\alpha:\Pi\to\Pi'$ carrying $J$ onto $J'$
induces $E_J\to E_{J'}$ by conjugation on automorphisms. Applying the
JSJ theorem, (B.JSJ.2) and (B.JSJ.3) proves the stated functorial
reconstruction. If a ramified extension of degree $d$ is specified,
its inertia image is the corresponding open subgroup and the same
construction gives the base-changed thickness $dm_e$, as it must.
No thickness independent of the normalization of the trait is asserted.
\end{proof}
